% English translation of chapter1.tex, lines 33--139.
% Source: Methods of Algebra, Volume 2, commit
% 9a5803ff77dd3257484cb177f851a73770a59dd3 (CC BY 4.0).
% stable-unit-id: o014.aljabr2.chapter1.subquotients

% segment-id: o014.li2.chapter1.subquotients.h001
\section{Subquotients}\label{sec:subquot}
\hypertarget{o014.aljabr2.chapter1.subquotients}{ }

% segment-id: o014.li2.chapter1.subquotients.p001
We begin by recalling the notions of monomorphism and epimorphism. Fix a
category $\mathcal C$.

% segment-id: o014.li2.chapter1.subquotients.l001
\begin{itemize}
	% segment-id: o014.li2.chapter1.subquotients.i001
	\item A morphism $f:X\to Y$ is called a \emph{monomorphism}, also written
	$f:X\hookrightarrow Y$, if left cancellation holds:
	\[ \forall T \in \Obj(\mathcal{C}), \; \forall g,h: T \to X , \quad fg=fh \iff g=h; \]
	% segment-id: o014.li2.chapter1.subquotients.i002
	\item A morphism $f:X\to Y$ is called an \emph{epimorphism}, also written
	$f:X\twoheadrightarrow Y$, if right cancellation holds:
	\[ \forall T \in \Obj(\mathcal{C}), \; \forall g,h: Y \to T, \quad gf=hf \iff g=h. \]
\end{itemize}

% segment-id: o014.li2.chapter1.subquotients.x001
\index{monomorphism}
\index{epimorphism}

% segment-id: o014.li2.chapter1.subquotients.p002
It follows immediately that $f:X\to Y$ is monic if and only if
$f_*:\Hom(T,X)\to\Hom(T,Y)$ is injective for every $T$. Likewise, $f$ is
epic if and only if $f^*:\Hom(Y,T)\to\Hom(X,T)$ is injective for every $T$.
Monicity and epicity are dual properties: $f$ is monic if and only if
$f^{\opp}$ is epic.

% segment-id: o014.li2.chapter1.subquotients.e001
\begin{example}
	If the equalizer $\Ker(f,g)$ of a pair of morphisms $f,g:X\to Y$ exists,
	its universal property shows that the canonical morphism
	$\iota:\Ker(f,g)\to X$ is monic. Dually, if the coequalizer
	$\Coker(f,g)$ exists, the canonical morphism $p:Y\to\Coker(f,g)$ is epic.
	Here is another simple observation: if $\mathcal C$ has a zero object,
	denoted by $0$, then $0\to X$ is monic and $X\to0$ is epic. The reader is
	invited to verify these details.
\end{example}

% segment-id: o014.li2.chapter1.subquotients.pr001
\begin{proposition}\label{prop:epi-mono-isomorphism}
	Consider morphisms $X\xrightarrow{a}Y\xrightarrow{b}Z$. Suppose that $b$
	is monic or that $a$ is epic. Then $ba:X\to Z$ is an isomorphism if and
	only if both $a$ and $b$ are isomorphisms.
\end{proposition}

% segment-id: o014.li2.chapter1.subquotients.pf001
\begin{proof}
	One direction is clear. Now suppose that $ba$ is an isomorphism and that
	$b$ is monic. Set $c:=a(ba)^{-1}:Z\to Y$; then $bc=\identity_Z$.
	Moreover, $b(cb)=(bc)b=b$, so left cancellation by $b$ gives
	$cb=\identity_Y$. Thus $b$ is an isomorphism, and consequently so is $a$.
	Passing to $\mathcal C^{\opp}$ gives the case in which $a$ is epic.
\end{proof}

% segment-id: o014.li2.chapter1.subquotients.p003
In categories commonly used in algebra, such as $\cate{Set}$, $\cate{Ab}$,
and $R\dcate{Mod}$, a morphism is monic or epic precisely when its underlying
map is injective or surjective, respectively. In these examples one can also
speak of subsets, subgroups, and so forth. This notion extends to arbitrary
categories, but one additional step is needed. We formulate it in the language
of partially ordered sets.

% segment-id: o014.li2.chapter1.subquotients.p004
Let $\mathcal C$ be a category and let $X$ be an object of $\mathcal C$.
Given a pair of monomorphisms $f_i:S_i\hookrightarrow X$, where $i=1,2$,
write $(S_1,f_1)\subset(S_2,f_2)$, or simply $S_1\subset S_2$, if there is
a morphism $g:S_1\to S_2$ such that $f_2g=f_1$. The monicity of $f_2$ implies
that $g$, if it exists, is unique and is itself monic. The relation $\subset$
defined in this way is not yet a partial order. If $S_1\subset S_2$ and
$S_2\subset S_1$, we say that $S_1$ and $S_2$ are equivalent. This is the
same as requiring a commutative diagram

% segment-id: o014.li2.chapter1.subquotients.g001
\[\begin{tikzcd}
	S_1 \arrow[r, "g_1"] \arrow[hookrightarrow, rd, "f_1"'] & S_2 \arrow[hookrightarrow, d, "f_2"] \arrow[r, "g_2"] & S_1 \arrow[hookrightarrow, ld, "f_1"] \\
	& X &
\end{tikzcd}\]

% segment-id: o014.li2.chapter1.subquotients.p005
Uniqueness gives $g_2g_1=\identity_{S_1}$ and, similarly,
$g_1g_2=\identity_{S_2}$. Thus $S_1$ and $S_2$ are equivalent if and only if
there is an isomorphism $g:S_1\rightiso S_2$ with $f_2g=f_1$; moreover, this
isomorphism is unique. This is, of course, a standard argument in algebra.

% segment-id: o014.li2.chapter1.subquotients.p006
Similarly, given a pair of epimorphisms $f_i:X\twoheadrightarrow Q_i$, where
$i=1,2$, if there is a morphism $g:Q_2\to Q_1$ such that $gf_2=f_1$, then
$g$ is unique and epic; in this case write $Q_1\twoheadleftarrow Q_2$. If
$Q_1\twoheadleftarrow Q_2$ and $Q_2\twoheadleftarrow Q_1$, we say that $Q_1$
and $Q_2$ are equivalent. Equivalently, there is a unique isomorphism
$g:Q_2\rightiso Q_1$ satisfying $gf_2=f_1$.

% segment-id: o014.li2.chapter1.subquotients.p007
The binary relation $\subset$ (or $\twoheadleftarrow$) induces a partial order
on the corresponding equivalence classes. We have borrowed the set-inclusion
symbol $\subset$ here. Although this may cause some ambiguity, it is consistent
with the notation for other algebraic substructures, such as subgroups,
submodules, and subspaces.

% segment-id: o014.li2.chapter1.subquotients.d001
\begin{definition}\label{def:sub-quot-obj}
	\index{subobject} \index[sym1]{SubX@$\mathrm{Sub}_X$}
	\index{quotient object} \index[sym1]{QuotX@$\mathrm{Quot}_X$}
	Let $\mathcal C$ be a category and let $X$ be an object of $\mathcal C$.
	An equivalence class, under the relation above, of monomorphisms
	$S\hookrightarrow X$ (respectively, of epimorphisms
	$X\twoheadrightarrow Q$) is called a \emph{subobject} (respectively, a
	\emph{quotient object}) of $X$. These equivalence classes form the
	partially ordered set $(\mathrm{Sub}_X,\subset)$ (respectively,
	$(\mathrm{Quot}_X,\twoheadleftarrow)$), in which
	$\identity_X:X\to X$ gives the unique maximal element.
\end{definition}

% segment-id: o014.li2.chapter1.subquotients.p008
Subobjects and quotient objects are dual: $(\mathrm{Sub}_X,\subset)$ as
defined in $\mathcal C$ is the same as
$(\mathrm{Quot}_X,\twoheadleftarrow)$ as defined in $\mathcal C^{\opp}$.
For convenience, this section deals mainly with subobjects. Because the
isomorphism $g$ occurring in the equivalence relation is unique, one may
choose a representative $f:S\hookrightarrow X$ of an equivalence class
without ambiguity, or abbreviate that representative simply as $S$. Recall
that $f_*:\Hom(\cdot,S)\to\Hom(\cdot,X)$ is injective.

% segment-id: o014.li2.chapter1.subquotients.lm001
\begin{lemma}\label{prop:subobject-order}
	For any two subobjects $S_1,S_2$ of $X$, one has
	\begin{equation*}
		S_1 \subset S_2 \iff \forall T \in \Obj(\mathcal{C}), \; \Hom(T, S_1) \subset \Hom(T, S_2),
	\end{equation*}
	where $\Hom(T,S_1)$ and $\Hom(T,S_2)$ on the right-hand side are
	understood as subsets of $\Hom(T,X)$.
\end{lemma}

% segment-id: o014.li2.chapter1.subquotients.pf002
\begin{proof}
	The implication $\implies$ is immediate. Conversely, take $T:=S_1$.
	Under $\Hom(S_1,S_1)\subset\Hom(S_1,X)$, the element
	$\identity_{S_1}$ corresponds to $f_1$. The hypothesis supplies
	$g\in\Hom(S_1,S_2)$ whose image in $\Hom(S_1,X)$ is $f_1$, that is,
	$f_2g=f_1$. This is the required morphism.
\end{proof}

% segment-id: o014.li2.chapter1.subquotients.p009
Suppose that $\mathcal C$ has a zero object and that $(X_i)_{i\in I}$ is a
family of objects. For every $(i,j)\in I\times I$, define
$\delta_{ij}\in\Hom(X_i,X_j)$ by $\delta_{ij}=\identity$ when $i=j$, and as
the zero morphism otherwise. If the product and coproduct of
$(X_i)_{i\in I}$ exist, then $(\delta_{ij})_{i,j}$ determines a canonical
morphism

% segment-id: o014.li2.chapter1.subquotients.q001
\begin{equation}\label{eqn:coprod-prod-delta}
	\delta: \coprod_{i \in I} X_i \to \prod_{i \in I} X_i.
\end{equation}

% segment-id: o014.li2.chapter1.subquotients.p010
For the additive categories to be reviewed shortly, $\delta$ is always an
isomorphism when $I$ is finite. This need not hold in an arbitrary category.

% segment-id: o014.li2.chapter1.subquotients.lm002
\begin{lemma}\label{prop:mono-pullback}
	Let $X\to Z$ be a monomorphism. Its pullback along $Y\to Z$, namely
	$X\dtimes{Z}Y\to Y$ (if it exists), is also a monomorphism. Dually, a
	pushout of an epimorphism is again an epimorphism.
\end{lemma}

% segment-id: o014.li2.chapter1.subquotients.pf003
\begin{proof}
	By duality, it suffices to treat the case of monomorphisms. Let
	$p_2:X\dtimes{Z}Y\to Y$ be the canonical morphism supplied by the fiber
	product. The question reduces to showing, for every object $T$, that
	% segment-id: o014.li2.chapter1.subquotients.g002
	\[\begin{tikzcd}
		\Hom(T,X) \dtimes{\Hom(T,Z)} \Hom(T,Y) & \\
		\Hom\left(T, X \dtimes{Z} Y \right) \arrow[leftrightarrow, u, "1:1", "\text{universal property}"'] \arrow[r, "{(p_2)_*}"'] & \Hom(T, Y)
	\end{tikzcd} \]
	is injective. Indeed, the monicity of $X\to Z$ makes
	$\Hom(T,X)\to\Hom(T,Z)$ injective, and the map from the upper left to the
	lower right is the evident projection from the fiber product.
\end{proof}

% segment-id: o014.li2.chapter1.subquotients.lm003
\begin{lemma}\label{prop:mono-product}
	Let $(f_i:X_i\to Z)_{i\in I}$ be a family of monomorphisms. Then the
	canonical morphism from $\prod_{i\in I}(X_i\to Z)$ (if it exists) to $Z$
	is also a monomorphism. Dually, given a family of epimorphisms
	$(g_i:Z\to X_i)_{i\in I}$, the canonical morphism from $Z$ to
	$\coprod_{i\in I}(Z\to X_i)$ (if it exists) is also an epimorphism.
\end{lemma}

% segment-id: o014.li2.chapter1.subquotients.pf004
\begin{proof}
	As in Lemma~\ref{prop:mono-pullback}, the question reduces to showing that
	the bottom row of
	% segment-id: o014.li2.chapter1.subquotients.g003
	\[\begin{tikzcd}
		\prod_{i \in I} \left( \Hom(T, X_i) \to \Hom(T, Z) \right) & \\
		\Hom\left(T, \prod_{i \in I} (X_i \to Z) \right) \arrow[leftrightarrow, u, "1:1", "\text{universal property}"'] \arrow[r] & \Hom(T, Z)
	\end{tikzcd}\]
	is injective. This follows from the hypothesis, since every map
	$\Hom(T,X_i)\to\Hom(T,Z)$ is injective.
\end{proof}

% segment-id: o014.li2.chapter1.subquotients.p011
We return to subobjects and quotient objects.

% segment-id: o014.li2.chapter1.subquotients.d002
\begin{definition}\label{def:subquotient}
	\index{subquotient}
	Let $\mathcal C$ be any category and let $X,Y$ be objects of $\mathcal C$.
	If $Y$ can be realized as a quotient object of a subobject of $X$, then
	$Y$ is called a \emph{subquotient} of $X$.
\end{definition}

% segment-id: o014.li2.chapter1.subquotients.p012
Considering a subobject of a quotient object amounts to discussing a
subquotient in $\mathcal C^{\opp}$. The next result shows that, for a certain
class of categories, the two approaches yield the same notion. This class
includes the abelian categories studied later; see
Corollary~\sourcecrossref{prop:Abel-cat-subquotient}{2.1.7}.

% segment-id: o014.li2.chapter1.subquotients.pr002
\begin{proposition}\label{prop:subquotient-equiv}
	Suppose that $\mathcal C$ has finite fiber products and finite fiber
	coproducts, that pullbacks of epimorphisms remain epimorphisms, and that
	pushouts of monomorphisms remain monomorphisms. Then a subquotient may
	equivalently be defined as a subobject of a quotient object.
\end{proposition}

% segment-id: o014.li2.chapter1.subquotients.pf005
\begin{proof}
	Consider a subquotient diagram $Y\twoheadleftarrow Z\hookrightarrow X$ in
	$\mathcal C$, and set $W:=X\dsqcup{Z}Y$. The morphism $X\to W$ is the
	pushout of $Z\to Y$, so it is epic by
	Lemma~\ref{prop:mono-pullback}; the morphism $Y\to W$ is the pushout of
	$Z\to X$, so it is monic by hypothesis. Thus
	$Y\hookrightarrow W\twoheadleftarrow X$, exhibiting a subquotient of $X$
	in $\mathcal C^{\opp}$. Since the assumptions are self-dual, a subquotient
	in $\mathcal C^{\opp}$ likewise yields a subquotient in $\mathcal C$.
\end{proof}

% segment-id: o014.li2.chapter1.subquotients.p013
Once the property ``quotient of a subobject = subobject of a quotient'' has
been established, it follows that a subquotient of a subquotient is again a
subquotient.

% segment-id: o014.li2.chapter1.subquotients.p014
The definition of a subquotient does not specify how it is to be realized as a
quotient object of a subobject. For example, in $\mathcal C=\cate{Ab}$,
$\Z/2\Z$ is a subquotient of $\Z/4\Z$: it can be realized either as a quotient
object of $\Z/4\Z$ or as the subobject $2\Z/4\Z$ of $\Z/4\Z$.
